% Evaluation-table design for the compact three-RQ organization.
% RQ4's ablation and robustness questions are incorporated into RQ1.
%
% Filled pilot values are reconstructed from the legacy 2026-07-24 workbooks:
%   48U_metrics_12MP_normal_0724.xlsx
%   48U_metrics_24MP_memory_0724.xlsx
% These workbooks predate evaluationSchemaVersion=2. They are used only to
% illustrate the table form and must be replaced by independent-session
% aggregates from the assigned policy arms.

\begin{table*}[t]
  \centering
  \caption{RQ1: Full-controller safety, feature delivery, and deadline headroom
  over burst prefixes. Safe is the percentage of sessions without a Capture
  Timeout or Draft watchdog through the prefix; Full is the percentage of
  planned shots completing both $M$ and $S$; Slack is P5/minimum deadline
  margin.}
  \label{tab:rq1_thermal_prefix}
  {\scriptsize
  \setlength{\tabcolsep}{2.1pt}%
  \renewcommand{\arraystretch}{1.10}%
  \begin{tabular}{@{}c c || c c c | c c c | c c c@{}}
    \toprule
    \multirow{2}{*}{\textbf{Device}} &
    \multirow{2}{*}{\makecell{\textbf{Start}\\\textbf{Lv}}} &
    \multicolumn{3}{c|}{\textbf{First 5 shots}} &
    \multicolumn{3}{c|}{\textbf{First 10 shots}} &
    \multicolumn{3}{c}{\textbf{First 30 shots}} \\
    \cmidrule(lr){3-5}\cmidrule(lr){6-8}\cmidrule(l){9-11}
    & & \makecell{Safe\\(\%)} & \makecell{Full\\(\%)} &
    \makecell{Slack\\P5/min (s)} &
    \makecell{Safe\\(\%)} & \makecell{Full\\(\%)} &
    \makecell{Slack\\P5/min (s)} &
    \makecell{Safe\\(\%)} & \makecell{Full\\(\%)} &
    \makecell{Slack\\P5/min (s)} \\
    \midrule
    \multicolumn{11}{c}{\textbf{(a) 12MP, normal memory}} \\
    \midrule
    \multirow{7}{*}{D1}
      & Lv0 & 100 & 100 & 5.95/5.90 & 100 & 100 & 6.01/5.90 &
        \multicolumn{3}{c}{--} \\
      & Lv1 & 100 & 100 & 6.23/6.23 & 100 & 100 & 6.22/6.22 &
        \multicolumn{3}{c}{--} \\
      & Lv2 & 100 & 100 & 5.90/5.88 & 100 & 100 & 5.91/5.88 &
        100 & 100 & 2.71/2.51 \\
      & Lv3 & 100 & 100 & 5.99/5.99 & 100 & 100 & 6.00/5.99 &
        100 & 100 & 4.08/3.78 \\
      & Lv4 & 100 & 100 & 1.52/1.48 & 100 & 100 & 0.61/0.51 &
        100 & 90.0 & 0.24/0.16 \\
      & Lv5 & 100 & 100 & 0.35/0.33 & 100 & 50.0 & 0.33/0.33 &
        100 & 43.3 & 0.40/0.33 \\
      & Lv6 & 100 & 60.0 & 0.15/0.15 & 100 & 30.0 & 0.15/0.15 &
        100 & 60.0 & 0.15/0.11 \\
    \midrule
    \multicolumn{11}{c}{\textbf{(b) 24MP, memory pressure}} \\
    \midrule
    \multirow{7}{*}{D1}
      & Lv0 & \multicolumn{3}{c|}{--} & \multicolumn{3}{c|}{--} &
        \multicolumn{3}{c}{--} \\
      & Lv1 & 100 & 100 & 3.38/3.32 & 100 & 100 & 2.05/1.93 &
        \multicolumn{3}{c}{--} \\
      & Lv2 & 100 & 0.0 & 1.57/1.52 & 100 & 30.0 & 1.64/1.52 &
        100 & 76.7 & 0.38/0.23 \\
      & Lv3 & 100 & 100 & 0.41/0.41 & 100 & 100 & 0.39/0.37 &
        100 & 66.7 & 0.36/0.18 \\
      & Lv4 & 100 & 0.0 & 1.99/1.99 & 100 & 0.0 & 1.99/1.99 &
        100 & 26.7 & 0.53/0.37 \\
      & Lv5 & 100 & 60.0 & 0.41/0.39 & 100 & 30.0 & 0.43/0.39 &
        100 & 30.0 & 0.52/0.39 \\
      & Lv6 & 100 & 0.0 & 3.15/3.14 & 100 & 30.0 & 3.01/2.90 &
        100 & 23.3 & 0.34/0.20 \\
    \bottomrule
  \end{tabular}%
  }
  \par\vspace{3pt}
  {\footnotesize Higher Safe, Full, and Slack values are better. The pilot
  cells above group one continuous legacy trace by observed level and are only
  layout examples. Confirmatory cells must aggregate independent sessions by
  \emph{starting} level, retain failed/truncated sessions in the planned-shot
  denominator, and repeat each seven-level block for D2 and D3 and for all four
  resolution--memory conditions.}
\end{table*}

\begin{table*}[t]
  \centering
  \caption{RQ2: Factual admission correctness and upper-bound calibration.}
  \label{tab:rq2_admission_correctness}
  {\small
  \setlength{\tabcolsep}{4.0pt}%
  \renewcommand{\arraystretch}{1.12}%
  \begin{tabular}{@{}c l r c c c c r c@{}}
    \toprule
    \textbf{Device} & \textbf{Condition} &
    \makecell{\textbf{Audited A/S}\\\textbf{(\#)}} &
    \makecell{\textbf{Correct}\\\textbf{admit} $\uparrow$} &
    \makecell{\textbf{Unsafe}\\\textbf{admit} $\downarrow$} &
    \makecell{\textbf{Unnecessary}\\\textbf{skip} $\downarrow$} &
    \makecell{\textbf{UB coverage}\\$\uparrow$} &
    \makecell{\textbf{P95 UB miss}\\\textbf{(ms)} $\downarrow$} &
    \makecell{\textbf{Admit audit}\\\textbf{coverage} $\uparrow$} \\
    \midrule
    D1 & 12MP normal & 1804/0 & 100.0\% & 0.0\% & -- &
      84.4\% & 734 & 99.2\% \\
    D1 & 24MP memory & 1554/0 & 100.0\% & 0.0\% & -- &
      83.1\% & 1242 & 98.9\% \\
    \midrule
    D2 & All $2{\times}2$ conditions & -- & -- & -- & -- & -- & -- & -- \\
    D3 & All $2{\times}2$ conditions & -- & -- & -- & -- & -- & -- & -- \\
    \bottomrule
  \end{tabular}%
  }
  \par\vspace{3pt}
  {\footnotesize A decision is factually safe when its observed remaining Draft
  wall time is no greater than the usable budget at the decision. Thus, a safe
  admit is correct and an unsafe admit is an error. UB coverage asks whether the
  predicted upper bound covers that factual wall cost and is independent of the
  later timeout outcome. A/S gives factual admit/skip audit counts. Normal
  execution does not reveal the cost of skipped work, so unnecessary-skip cells
  are intentionally blank; objective values require a node-level randomized
  forced-execution audit. Report session-clustered intervals rather than treating
  decisions as independent samples.}
\end{table*}

\begin{table*}[t]
  \centering
  \caption{RQ3: Full-controller pacing profile over thermal burst prefixes.}
  \label{tab:rq3_thermal_prefix}
  {\scriptsize
  \setlength{\tabcolsep}{1.9pt}%
  \renewcommand{\arraystretch}{1.10}%
  \begin{tabular}{@{}c c || c c c | c c c | c c c@{}}
    \toprule
    \multirow{2}{*}{\textbf{Device}} &
    \multirow{2}{*}{\makecell{\textbf{Start}\\\textbf{Lv}}} &
    \multicolumn{3}{c|}{\textbf{First 5 shots}} &
    \multicolumn{3}{c|}{\textbf{First 10 shots}} &
    \multicolumn{3}{c}{\textbf{First 30 shots}} \\
    \cmidrule(lr){3-5}\cmidrule(lr){6-8}\cmidrule(l){9-11}
    & & \makecell{Paced\\(\%)} & \makecell{Delay\\mean/P95} &
    \makecell{S2S\\P95} &
    \makecell{Paced\\(\%)} & \makecell{Delay\\mean/P95} &
    \makecell{S2S\\P95} &
    \makecell{Paced\\(\%)} & \makecell{Delay\\mean/P95} &
    \makecell{S2S\\P95} \\
    \midrule
    \multicolumn{11}{c}{\textbf{(a) 12MP, normal memory}} \\
    \midrule
    \multirow{7}{*}{D1}
      & Lv0 & 0.0 & 0/0 & 700 & 0.0 & 0/0 & 683 &
        \multicolumn{3}{c}{--} \\
      & Lv1 & -- & -- & 713 & 0.0 & 0/0 & 678 &
        \multicolumn{3}{c}{--} \\
      & Lv2 & 0.0 & 0/0 & 574 & 0.0 & 0/0 & 572 &
        0.0 & 0/0 & 589 \\
      & Lv3 & 0.0 & 0/0 & 708 & 0.0 & 0/0 & 683 &
        0.0 & 0/0 & 654 \\
      & Lv4 & 0.0 & 0/0 & 657 & 0.0 & 0/0 & 736 &
        27.6 & 54/243 & 851 \\
      & Lv5 & 0.0 & 0/0 & 850 & 50.0 & 153/346 & 1135 &
        20.7 & 72/346 & 1139 \\
      & Lv6 & 100 & 359/402 & 1113 & 90.0 & 339/445 & 1215 &
        44.8 & 163/456 & 1139 \\
    \midrule
    \multicolumn{11}{c}{\textbf{(b) 24MP, memory pressure}} \\
    \midrule
    \multirow{7}{*}{D1}
      & Lv0 & \multicolumn{3}{c|}{--} & \multicolumn{3}{c|}{--} &
        \multicolumn{3}{c}{--} \\
      & Lv1 & 0.0 & 0/0 & 609 & 0.0 & 0/0 & 623 &
        \multicolumn{3}{c}{--} \\
      & Lv2 & 100 & 142/229 & 894 & 77.8 & 136/283 & 896 &
        31.0 & 72/283 & 747 \\
      & Lv3 & 100 & 182/405 & 1181 & 100 & 171/373 & 1123 &
        44.8 & 165/983 & 1691 \\
      & Lv4 & 0.0 & 0/0 & 627 & 0.0 & 0/0 & 891 &
        13.8 & 109/890 & 1520 \\
      & Lv5 & 40.0 & 233/738 & 1061 & 50.0 & 233/691 & 1639 &
        17.2 & 80/445 & 1326 \\
      & Lv6 & 0.0 & 0/0 & 867 & 0.0 & 0/0 & 868 &
        13.8 & 168/1082 & 1558 \\
    \bottomrule
  \end{tabular}%
  }
  \par\vspace{3pt}
  {\footnotesize Delay and shot-to-shot (S2S) values are in ms; delay statistics
  include zero-delay decisions. This table shows when and how strongly the Full
  controller acts, not the causal benefit of pacing. That benefit is identified
  by Full$-$Admission-only and Pacing-only$-$Always-on in
  Table~\ref{tab:rq1_ablation}. The pilot-level caveat and D2/D3 expansion are
  the same as in Table~\ref{tab:rq1_thermal_prefix}.}
\end{table*}

\begin{table*}[t]
  \centering
  \caption{RQ1 ablation and robustness scorecard.}
  \label{tab:rq1_ablation}
  {\small
  \setlength{\tabcolsep}{4.0pt}%
  \renewcommand{\arraystretch}{1.08}%
  \begin{tabular}{@{}c l c c c c c c@{}}
    \toprule
    \textbf{Device} & \textbf{Assigned policy} &
    \makecell{\textbf{Safe@30}\\(\%) $\uparrow$} &
    \makecell{\textbf{Full@30}\\(\%) $\uparrow$} &
    \makecell{\textbf{$S$@30}\\(\%) $\uparrow$} &
    \makecell{\textbf{Slack P5}\\(s) $\uparrow$} &
    \makecell{\textbf{S2S P95}\\(ms) $\downarrow$} &
    \makecell{\textbf{30-shot span}\\(s) $\downarrow$} \\
    \midrule
    \multirow{6}{*}{D1}
      & Always-on & -- & -- & -- & -- & -- & -- \\
      & Static-L4 & -- & -- & -- & -- & -- & -- \\
      & Admission-only & -- & -- & -- & -- & -- & -- \\
      & Pacing-only & -- & -- & -- & -- & -- & -- \\
    \rowcolor{bgGray}
      & \textbf{Full} & \textbf{100/100} & \textbf{100/66.7} &
        \textbf{100/100} & \textbf{6.17/0.45} &
        \textbf{683/738} & \textbf{16.7/14.9} \\
      & Mandatory-only (diagnostic) & -- & -- & -- & -- & -- & -- \\
    \midrule
    \multirow{6}{*}{D2}
      & Always-on & -- & -- & -- & -- & -- & -- \\
      & Static-L4 & -- & -- & -- & -- & -- & -- \\
      & Admission-only & -- & -- & -- & -- & -- & -- \\
      & Pacing-only & -- & -- & -- & -- & -- & -- \\
    \rowcolor{bgGray}
      & \textbf{Full} & -- & -- & -- & -- & -- & -- \\
      & Mandatory-only (diagnostic) & -- & -- & -- & -- & -- & -- \\
    \midrule
    \multirow{6}{*}{D3}
      & Always-on & -- & -- & -- & -- & -- & -- \\
      & Static-L4 & -- & -- & -- & -- & -- & -- \\
      & Admission-only & -- & -- & -- & -- & -- & -- \\
      & Pacing-only & -- & -- & -- & -- & -- & -- \\
    \rowcolor{bgGray}
      & \textbf{Full} & -- & -- & -- & -- & -- & -- \\
      & Mandatory-only (diagnostic) & -- & -- & -- & -- & -- & -- \\
    \bottomrule
  \end{tabular}%
  }
  \par\vspace{3pt}
  {\footnotesize Static-L4 is the primary production baseline. Interpret the
  scorecard lexicographically: establish safety first, then compare feature
  yield among similarly safe policies, and finally compare responsiveness.
  D1 Full entries are illustrative 12MP-normal/24MP-memory pilot pairs; the
  other cells require separately assigned on-device arm runs. Final main-table
  entries should be session-balanced macro-estimates across the complete
  $2{\times}2$ resolution--memory matrix, with condition-specific results in
  the appendix.}
\end{table*}

\begin{table*}[t]
  \centering
  \caption{Pilot stress cases motivating matched preventive-rescue analysis.}
  \label{tab:preventive_intervention}
  {\small
  \setlength{\tabcolsep}{5.0pt}%
  \renewcommand{\arraystretch}{1.10}%
  \begin{tabular}{@{}c l r r r r c c@{}}
    \toprule
    \textbf{Device} & \textbf{Condition} & \textbf{Full shot} &
    \makecell{\textbf{Pacing}\\\textbf{delay (ms)}} &
    \makecell{\textbf{Queue wait}\\\textbf{(ms)}} &
    \makecell{\textbf{Slack}\\\textbf{(ms)}} &
    \textbf{Output} & \textbf{Outcome} \\
    \midrule
    D1 & 12MP normal & 293 & 358 & 5067 & $+75$ & $M{+}S$ & Safe \\
    D1 & 24MP memory & 48 & 723 & 5100 & $+229$ & $M{+}S$ & Safe \\
    \bottomrule
  \end{tabular}%
  }
  \par\vspace{3pt}
  {\footnotesize These rows show that the requested delay was absorbed while
  the Draft queue remained busy and that the stressed capture still delivered
  $M{+}S$ with positive margin. They are mechanism candidates, not causal
  rescues, because the pilot lacks a matched no-control timeout. The final case
  table should align each baseline's first timeout $k_0$ with Full shots
  $k_0-2$, $k_0-1$, and $k_0$, report the first skip/delay and lead shots
  $k_0-k_{\mathrm{first\ action}}$, and call a pair rescued only when Full is
  failure-free through $k_0$ and completes $M{+}S$ at $k_0$. The aggregate
  result is the rescued-pair rate over all matched baseline failures.}
\end{table*}
